\usepackage{amssymb,amsmath,mathabx, amscd}
\usepackage[pdftex]{graphicx}
\usepackage{enumitem}
\usepackage{tikz, tikz-cd}
\usepackage[margin=.8in]{geometry}
\usepackage{xcolor}
\usepackage{hyperref}
\usepackage{mathrsfs, colonequals}
\usepackage{fancyhdr}

\newcommand{\nn}{\mathbb{N}}
\newcommand{\cc}{\mathbb{C}}
\newcommand{\ff}{\mathbb{F}}
\newcommand{\qq}{\mathbb{Q}}
\newcommand{\rr}{\mathbb{R}}
\newcommand{\zz}{\mathbb{Z}}
\newcommand{\bb}{\mathbb{B}}
\newcommand{\pp}{\mathbb{P}}

\newcommand{\B}{\mathcal{B}}
\renewcommand{\P}{\mathcal{P}}

\newtheorem{thm}{Theorem}[section]
\newtheorem{ithm}{Theorem}
\newtheorem{icor}[ithm]{Corollary}
\newtheorem{lem}[thm]{Lemma}
\newtheorem{conj}[thm]{Conjecture}
\newtheorem{prop}[thm]{Proposition}
\newtheorem{cor}[thm]{Corollary}
\newtheorem{claim}[thm]{Claim}

\theoremstyle{definition}
\newtheorem{defin}[thm]{Definition}
\newtheorem{example}[thm]{Example}
\newtheorem{assumption}[thm]{Assumption}
\newtheorem*{question}{Main Question}

\theoremstyle{remark}
\newtheorem{rmk}{Remark}[section]

\newtheorem{exercise}{Exercise}[section]

\newcommand{\nul}{\operatorname{null}}
\newcommand{\range}{\operatorname{range}}
\newcommand{\Span}{\operatorname{span}}
\newcommand{\Real}{\operatorname{Re}}
\renewcommand{\implies}{\Longrightarrow}
\renewcommand{\iff}{\Leftrightarrow}
\newcommand{\diag}{\operatorname{diag}}
\newcommand{\defi}[1]{{\color{violet}\textsf{#1}}} 
\DeclareMathOperator{\tr}{tr}
\newcommand{\interior}{\operatorname{Int}}
\newcommand{\closure}{\operatorname{Cl}}

\newcommand{\comment}[1]{\textcolor{magenta}{#1}}
\def\multiset#1#2{\ensuremath{\left(\kern-.2em\left(\genfrac{}{}{0pt}{}{#1}{#2}\right)\kern-.2em\right)}}
\newenvironment{ans}
{\color{teal}\par\noindent\textbf{Answer:}\noindent\begin{itshape}}
{\end{itshape}\par}

